%% ma: L10-B02-0044
%% lop: 10
%% chuong: 1
%% bai: 2
%% dang: 10.2.7
%% loai: TLN
%% muc: VD
%% dap_an: "2,67"
%% nguon: "Ảnh giáo viên gửi 10/10/2026 (ThS. Nguyễn Hoàng Việt, Chương 2 Tập hợp), Đề 2 (đề thứ hai, không ghi tiêu đề), Phần 3 Câu 22"
%% kien_thuc: "Bài 2 (tập hợp và các phép toán trên tập hợp)"
%% trang_thai: cho-duyet
%% ly_do: "Dạng toán mới đề xuất 10.2.7: tìm tham số để các tập hợp thỏa quan hệ (giao khác rỗng, tập con, rời nhau)"
%% ghi_chu: "Hiểu A khác rỗng (m < 14/3) theo dạng m thuộc [a;b) của đề. Nếu cho phép A rỗng thì mọi m ≥ 14/3 cũng thỏa, khi đó tập m không có dạng [a;b)"
\begin{ex}
Cho $A=\left[m-3;\dfrac{m+2}{4}\right)$, $B=(-\infty;-1)\cup[2;+\infty)$. Khi $m\in[a;b)$ thì $A\cap B=\varnothing$. Khi đó $b-a$ bằng (kết quả làm tròn đến hàng phần trăm).
\shortans{2,67}
\loigiai{Để $A$ là nửa khoảng khác rỗng: $m-3<\dfrac{m+2}{4}\Leftrightarrow m<\dfrac{14}{3}$. $A\cap B=\varnothing\Leftrightarrow A\subset[-1;2)\Leftrightarrow m-3\ge-1$ và $\dfrac{m+2}{4}\le2\Leftrightarrow2\le m\le6$. Kết hợp với $m<\dfrac{14}3$: $m\in\left[2;\dfrac{14}{3}\right)$, nên $a=2$, $b=\dfrac{14}3$ và $b-a=\dfrac83\approx2{,}67$.}
\end{ex}
